\begin{minipage}{0.55\linewidth}
Une tige horizontale uniforme de \qty{\v{M}}{kg} et de \qty{\v{L}}{m} est fixée à un mur par une charnière. Une enseigne de \qty{\v{m}}{kg} est suspendue à son extrémité libre. Un câble relie cette extrémité au mur et fait un angle de \qty{\v{th}}{\degree} avec la tige.
\end{minipage}\hfill
\begin{minipage}{0.4\linewidth}
\centering
\begin{tikzpicture}[scale=0.8]
  \fill[pattern=north east lines] (-0.3,-1.6) rectangle (0,2); \draw[thick] (0,-1.6) -- (0,2);
  \draw[line width=2.5pt] (0,0) -- (4,0); \fill (0,0) circle (2.5pt);
  \draw (0,1.8) -- (4,0);
  \draw (4,0) -- (4,-0.5); \draw[fill=gray!20] (3.3,-0.5) rectangle (4.7,-1.3);
  \draw (3.1,0) arc (180:156:0.9); \node at (2.8,0.25) {$\theta$};
\end{tikzpicture}
\end{minipage}

\begin{enumerate}[label=\alph*)]
    \item Déterminez la tension dans le câble.
    \item Déterminez la composante horizontale de la force exercée par la charnière sur la tige.
    \item Déterminez la composante verticale de la force exercée par la charnière sur la tige.
\end{enumerate}
